\documentclass[a4paper]{article}
\usepackage{amsfonts}
\usepackage{amsmath}
\usepackage{amssymb}
\usepackage{graphicx}     
\usepackage{cancel}

% Command definities van frequent gebruikte commandos
\newcommand{\N}{\mathbb{N}}
\newcommand{\Q}{\mathbb{Q}}
\newcommand{\R}{\mathbb{R}}
\newcommand{\C}{\mathbb{C}}
\newcommand{\Bgsin}{\operatorname{Bgsin}}
\newcommand{\Bgcos}{\operatorname{Bgcos}}
\newcommand{\Bgtan}{\operatorname{Bgtan}}
\newcommand{\cis}{\operatorname{cis}}
\newcommand{\Limit}{\lim\limits_}
\newcommand{\Int}{\displaystyle\int}
\newcommand{\Dint}{\displaystyle\int\limits}

\begin{document}
  
\noindent \large Auteur: Wout Vanderborght\\
\noindent \large Studierichting: Informatica\\
\noindent \large Oefeningenreeks 5A nr. 6.10\\

\medskip

\normalsize

$\Dint^9_2 \dfrac{3x}{\sqrt[3]{x-1}}dx $ \\

Stel $ t = x - 1 \Leftrightarrow x = t + 1 \Rightarrow dx = dt $. \\

$ \Rightarrow \Dint^9_2 \dfrac{3t+3}{t^{\tfrac{1}{3}}} = 3\Dint^9_2 t^{\tfrac{2}{3}} + 3\Dint^9_2 t^{\tfrac{-1}{3}} $ \\

$ = 3\left[ \dfrac{3t^{\tfrac{5}{3}}}{5} + \dfrac{3t^{\tfrac{2}{3}}}{2} \right]^9_2 = \left[ \dfrac{6\sqrt[3]{(x-1)^5} + 15\sqrt[3]{(x-1)^2}}{10} \right]^9_2 $ \\

$ = 3\left( \dfrac{6\cdot32 + 15\cdot4}{10} - \dfrac{6 + 15}{10} \right) = \dfrac{693}{10} $


\begin{thebibliography}{999}
%\bibitem{a} Referentie
\end{thebibliography}
\end{document} 
